\section{引言：在动态环境中学习}

本讲是 MIT 6.S191（Introduction to Deep Learning）系列课程的第五讲，主题为 \textbf{Deep Reinforcement Learning}（深度强化学习）。在前四讲中，课程已经覆盖了深度学习的核心架构（全连接网络、RNN、CNN）以及两类学习范式（Supervised Learning 和 Unsupervised Learning）。本讲将引入第三类学习范式------\textbf{Reinforcement Learning}（强化学习），并探讨如何将深度学习技术与强化学习相结合，使模型能够在\textbf{动态环境}中通过与环境的交互来学习最优行为策略。

\begin{figure}[H]
\centering
\includegraphics[width=0.85\textwidth]{slides-images/slide-02.jpg}
\caption{在动态环境中学习：强化学习让智能体通过与环境的交互来学习决策\protect\footnotemark}
\end{figure}
\footnotetext{Slides 第2页。}

深度强化学习的核心思想与之前的学习范式有本质区别：它不再依赖于预先收集好的静态数据集，而是让模型（Agent）在一个\textbf{动态环境}（Environment）中不断地进行交互------采取行动（Action）、观察环境变化（Observation）、获取奖励信号（Reward）------从而学会如何做出最优决策。这种学习方式在机器人控制、游戏策略、自动驾驶等领域具有广泛的应用价值。

\begin{knowledgebox}{为什么需要强化学习？}
在现实世界中，许多任务无法简单地转化为"给定输入预测输出"的监督学习问题。例如，让机器人学会行走、让 AI 学会下棋，都需要智能体在一系列时间步上做出连续决策，并且当前决策会影响未来的状态和奖励。这正是强化学习的核心应用场景。
\end{knowledgebox}

\begin{figure}[H]
\centering
\includegraphics[width=0.85\textwidth]{slides-images/slide-03.jpg}
\caption{强化学习的典型应用领域：机器人技术与游戏策略\protect\footnotemark}
\end{figure}
\footnotetext{Slides 第3页。}


